\begin{problem}{분탕의 신}{standard input}{standard output}{1 second}{256 megabytes}

이 문제는 인터랙티브 문제이다.

분탕의 신 문성이는 이번에도 가만히 있지 않았다.

정수 $N$이 주어진다.
그는 길이 $N$의 이진 수열 $A_1, A_2, \dots, A_N$을 만들어 두고, SCCC 부원들이 쉽게 알아내지 못하도록 온갖 장난을 쳐 두었다.
각 원소 $A_i$는 $0$ 또는 $1$이다.

당신의 목표는 채점기와 상호작용하여 \textbf{현재 수열} $A$를 정확히 알아내는 것이다.

당신은 채점기에게 최대 $1333$번까지 질문할 수 있다.
한 번의 질문은 다음 형식으로 한다.
\begin{center}
\texttt{? } $i$
\end{center}

이 질문을 하면 채점기는 현재 수열의 $i$번째 원소 $A_i$의 값을 반환한다.

하지만 문성이는 호락호락하지 않다.
채점기는 \textbf{10번째, 20번째, 30번째, \dots} 질문에 답한 직후마다 현재 수열에 대해 다음 네 가지 연산 중 하나를 임의로 적용한다.
\begin{itemize}
    \item 아무 연산도 하지 않음
    \item 전체 수열을 flip
    \item 전체 수열을 reverse
    \item 전체 수열을 flip한 뒤 reverse
\end{itemize}

flip과 reverse의 정의는 아래 노트를 참고하라.

현재 수열을 정확히 알아냈다면, 다음 형식으로 답을 출력해야 한다.
\begin{center}
\texttt{! } $A_1$ $A_2$ $\dots$ $A_N$
\end{center}

이때 출력한 수열은 \textbf{출력하는 순간의 현재 수열}과 정확히 일치해야 한다.

\InputFile
상호작용이 시작되면 첫째 줄에 정수 $N$이 주어진다.

\begin{itemize}
    \item $1 \le N \le 1000$
\end{itemize}

\OutputFile
현재 수열을 알아냈다면 다음 형식으로 한 줄을 출력한다.

\begin{itemize}
    \item \texttt{! }$A_1$ $A_2$ $\dots$ $A_N$
\end{itemize}

여기서 $A_1, A_2, \dots, A_N$은 공백으로 구분하여 출력해야 한다.
즉, \texttt{! 0 1 0 1}과 같은 형식으로 출력해야 하며, \texttt{! 0101}과 같이 붙여서 출력해서는 안 된다.

질문은 최대 1333번 까지 가능하다.

\Interaction
질문 \texttt{? i}를 출력하면, 채점기는 현재 수열의 $i$번째 값인 $0$ 또는 $1$을 한 줄에 출력한다.

참가자는 최종적으로 \texttt{! }$A_1$ $A_2$ $\dots$ $A_N$을 출력해야 하며, 이 수열은 출력하는 순간의 현재 수열과 정확히 같아야 한다.
여기서 $A_1, A_2, \dots, A_N$은 공백으로 구분하여 출력해야 한다.

출력을 즉시 채점 프로그램에 전달하기 위해, 질문이나 정답을 출력한 뒤에는 반드시 출력 버퍼를 비워야 한다.

언어별로 출력 버퍼를 비우는 방법의 예시는 다음과 같다.

\begin{itemize}
    \item \texttt{C}: \texttt{fflush(stdout);}
    \item \texttt{C++}: \texttt{std::cout.flush();}
    \item \texttt{Java}: \texttt{System.out.flush();}
    \item \texttt{Python}: \texttt{stdout.flush()}
\end{itemize}

그 외의 언어를 사용하는 경우에는 해당 언어의 문서를 참고하라.

정답을 출력한 뒤 프로그램은 즉시 종료해야 한다.

\Example

\begin{example}
\exmpfile{example.01}{example.01.a}%
\end{example}

\Note
\begin{itemize}
    \item \texttt{flip}은 수열의 모든 원소를 뒤집는 연산이다. 즉, 각 $i$에 대해 현재 값 $A_i$를 $1-A_i$로 바꾼다.
    \item \texttt{reverse}는 수열의 순서를 뒤집는 연산이다. 즉, 연산 후의 수열 $B$는 각 $i$에 대해 $B_i = A_{N+1-i}$를 만족한다.
    \item \texttt{flip 후 reverse}는 먼저 모든 원소를 뒤집은 뒤, 그 결과 수열의 순서를 뒤집는 연산이다. 즉, 연산 전 수열이 $A$이고 연산 후 수열이 $B$라면 각 $i$에 대해 $B_i = 1 - A_{N+1-i}$이다.
    \item 채점기의 변형은 오직 $10$번째, $20$번째, $30$번째, $\dots$ 질문에 답한 직후에만 일어날 수 있다. 최종 답안 출력 \texttt{! }$A_1$ $A_2$ $\dots$ $A_N$ 은 질문으로 세지지 않는다. 따라서 예를 들어 $9$번째 질문까지 한 뒤 곧바로 정답을 출력한다면, 그 사이에 수열이 추가로 변하는 일은 없다.
    \item 채점기가 어떤 연산을 적용했는지는 참가자에게 알려지지 않는다.
    \item 최종적으로 출력해야 하는 것은 초기 수열이 아니라, 정답을 출력하는 시점의 현재 수열이다.
\end{itemize}

\end{problem}

